
\section{Discussion}\label{sec:discussion}

For a neuron whose spike count in a window is Poisson, Barlow's question has an exact answer. The most informative rate map is a
staircase, its steps sit at the quantiles of the stimulus distribution, and its values are the mass points of the capacity-achieving
input of the counting channel (Proposition~\ref{prop:L1}). The stimulus distribution decides only where the steps sit. How many steps
there are depends on the peak count and the spontaneous count alone, and Theorem~\ref{thm:C} proves that the number grows at least like
$0.056\,A^{2/3}$, faster than the $\sqrt A$ rates that a count of resolvable levels would allow.

Laughlin's rule is the limit of this answer when noise is negligible: the response follows the distribution function of the stimulus.
With Poisson noise the equalization takes place in the square root of the rate, and the staircase approaches that curve as the peak
count grows. The computations of Nikitin and colleagues found staircases numerically~\cite{Nikitin2009}; the proposition gives them
for every continuous stimulus law, and the theorem fixes how fast their number grows. Kostal and Shinomoto found two levels for a
neuron observed in continuous time~\cite{KostalShinomoto2016}; counting in a window is what makes the number of levels grow.

Pooled over fibres, the auditory-nerve fibres favour the predicted shape, and they use few rates, close to the optimal number. Of the registered
tests on the fibres, the tests of shape held, the tests of discreteness under Poisson counts failed, and the test that the counts are
as variable as Poisson counts failed. Of the five tests of discreteness with regular counts, the two that count significant
differences held and the three others failed, including both likelihood tests, which have the most power. Their rate-level functions favour the asymmetric shape predicted for a bell-shaped distribution of sound levels over the
symmetric sigmoid, in young and in aged animals, and over a curve of the opposite asymmetry. A curve of the same direction with a
stronger asymmetry fits them better, so the data support the direction of the asymmetry more than the exact form. Under Poisson counts about half of them show more distinct rates than the staircase
allows, mostly one more, and neither drift across presentations nor a longer window explains the excess. Their counts, however, have
about a third of the variance of Poisson counts with the same mean, and a less noisy count supports more rates. For a count matched to the mean dispersion
of each fibre, $82.3\%$ of the fibres stay within the staircase, and $88.3\%$ do for the most regular binomial count
with the same range. For a count whose dispersion follows the measured fall of the Fano factor with the rate, $77.0\%$ do, below the
registered $80\%$. The agreement is weakest in the least regular third of the fibres, where $65.9\%$ to $84.1\%$ stay within it,
depending on the count model. The number of steps depends on
the 50~ms window through the peak count and on the spontaneous count. The distribution of sound levels sets only where the steps sit. For analyses of efficient coding with spike counts, the noise model decides the predicted number
of rates. Under the Poisson model about half of the fibres use more rates than the optimum. Under models with their measured
regularity most of them are consistent with the optimal staircase, from $77.0\%$ to $88.3\%$ of the fibres depending on how the
regularity is modelled. A likelihood test with more power, on the range from each fibre's spontaneous count, finds $49.4\%$ consistent with the optimal
number of steps and $86.8\%$ with one step more. The fibres thus use few rates; about half need more than the optimum, most of them one more.

The regularity of the counts has a known physiological source. An auditory-nerve fibre cannot fire again within an absolute refractory
period of about 0.6~ms~\cite{Heeringa2024}. Refractoriness makes the spike train of a neuron firing near its maximal rate more regular
than a Poisson process with the same rate~\cite{BerryMeister1998}. The binomial count of the tests above, at most one spike in each
short bin, is the simplest model of this effect. The theory treats one fibre and its count in one window, and three things lie outside
it. First, each sound level is coded by many fibres. Dynamic range adaptation is found in fibres of all spontaneous rates, and it
improves the precision of level coding by the fibre population~\cite{Wen2009}. Fibres whose steps sit at different sound levels can
together resolve more levels than any one of them; whether the population code is optimal is not tested here. Second, the count
discards spike timing, which a reader of the full spike train would use, as in the continuous-time result above. Third, spikes cost
energy. Under an energy budget at a peak count of 3, a graded code gains over the best on-off code at most $0.8\%$ of the capacity
at that budget (S1 Text).

Two limits of the present tests are stated where they arise. The shape test used the form of the optimal map for a small spontaneous
count. The exact form for any spontaneous count, given in Section~\ref{sec:laughlin}, was tested as well and describes the fibres
worse, as if the spontaneous count were an offset outside the range over which the information is optimized; why this is so is open. The
conjectured law for the exact number of steps held to within 0.15 in the peak count up to 47 steps and failed at 67. There it is
applied below the range on which its chain fit was made, and three further predictions of it are recorded and not yet computed. This
failure does not touch the tests on fibres: they use the number of steps computed directly for each fibre, at peak counts far below
those where the law failed. On six further public datasets from other species and areas, whose counts are close to Poisson counts, an
earlier test of a different design did not find discrete levels (S1 Text). That test lies outside the conditions of the theory in two
further ways. The window in which cortical neurons are read is not known, and the number of steps grows with the window. The stimuli
were a few conditions chosen by the experimenter, and three of the datasets come from motor and frontal areas, which do not encode
a stimulus distribution in Barlow's sense.
